\chapter{Introdução}

Na Engenharia de Software contemporânea, a Garantia da Qualidade (\textit{Quality Assurance} - QA) e a aplicação de processos metodológicos consolidados são pilares estratégicos indispensáveis para o sucesso de qualquer sistema de informação \cite{vasconcelos2006}. À medida que as aplicações digitais se tornam complexas e integradas às atividades organizacionais, a ocorrência de falhas ou comportamentos imprevistos passa a acarretar prejuízos financeiros substanciais, vulnerabilidades de segurança e falhas operacionais críticas \cite{koscianski2007}. Nesse cenário, as atividades de teste de software surgem como o principal mecanismo de verificação e validação, encarregadas de assegurar, de forma sistemática, que o produto final atenda rigorosamente aos requisitos funcionais e não funcionais especificados \cite{delamaro2016}.

Apesar da relevância das práticas de QA, a sua condução impõe desafios operacionais significativos às equipes de desenvolvimento. O ciclo de testes engloba tarefas complexas e interdependentes, tais como o planejamento do escopo de validação, a elaboração de casos de teste, a execução de cenários e o registro minucioso de defeitos. Para que este fluxo gere melhorias reais na maturidade do código, é essencial haver padronização terminológica e centralização das informações \cite{tebes2022}. Sistemas dedicados à gestão de testes são fundamentais para unificar artefatos e apoiar a organização do processo \cite{kokubum2025}. Contudo, em pequenas equipes ou ambientes acadêmicos, é frequente o recurso a abordagens \textit{ad-hoc} e ferramentas genéricas como planilhas eletrônicas ou documentos de texto compartilhados em nuvem para gerir as rotinas de teste. Esta fragmentação resulta em problemas crônicos, tais como a perda de histórico de execução, a desorganização de dados e a extrema dificuldade na extração de métricas de qualidade confiáveis \cite{medeiros2022}.

No contexto do Instituto Federal do Espírito Santo (IFES) Campus Alegre, essa realidade reflete-se diretamente nos projetos de software desenvolvidos no âmbito acadêmico, de pesquisa e de extensão. Atualmente, o gerenciamento dos ciclos de teste nesses projetos ocorre de maneira descentralizada e não padronizada, recorrendo frequentemente a documentos avulsos na nuvem, como a plataforma Google Docs. Essa falta de centralização dispersa o conhecimento técnico produzido, dificulta o acompanhamento do progresso dos projetos por parte dos professores orientadores e impede as equipes de extrair relatórios analíticos consolidados sobre a estabilidade do software desenvolvido. Torna-se necessária, portanto, a introdução de uma plataforma dedicada que centralize a gestão de QA e promova as boas práticas de Engenharia de Software na instituição.

Diante dessa necessidade prática, surge a proposta do ``Sistema Web para Gerenciamento de Casos de Teste e Rastreamento de Defeitos de Software: Test-Manager''. Trata-se de uma aplicação web delimitada para as equipes de desenvolvimento e QA do IFES Campus Alegre, cujo escopo abrange o cadastro de projetos, a criação e execução de casos de teste, o registro de defeitos e o controle hierárquico de acessos (Coordenador, Desenvolvedor e Tester). O sistema foi desenvolvido utilizando a linguagem PHP integrada ao banco de dados MySQL, um arranjo tecnológico consagrado pela robustez no desenvolvimento web \cite{welling2003}, adotando o framework Laravel sob o padrão de arquitetura MVC (\textit{Model-View-Controller}) \cite{subecz2021}. Essa direção tecnológica atende a uma solicitação direta do Professor Pedro, visando garantir a clareza estrutural do código e facilitar a manutenção e continuidade futura do projeto.

Dessa forma, este trabalho tem como objetivo geral desenvolver o sistema web Test-Manager para centralizar e aprimorar o controle de testes e o registro de defeitos nos projetos de software do IFES Campus Alegre, promovendo maior organização, preservação de histórico e governança nos ciclos de qualidade institucionais.

Este trabalho está estruturado em sete capítulos que detalham o desenvolvimento e a validação da pesquisa. Além desta Introdução --- que apresenta a Justificativa da proposta, delimita o Problema de Pesquisa por meio de uma pergunta norteadora e estabelece os Objetivos Geral e Específicos ---, o Capítulo 2 apresenta o Referencial Teórico sobre os conceitos essenciais da área, as tecnologias utilizadas e a análise comparativa de seis trabalhos correlatos da literatura. O Capítulo 3 descreve a Metodologia científica, o processo de desenvolvimento iterativo e incremental e o instrumento de validação baseado no Modelo TAM. O Capítulo 4 detalha o Projeto do Sistema Test-Manager, abrangendo sua arquitetura, os requisitos funcionais e não funcionais e a modelagem estrutural e comportamental em UML e no banco de dados. Em seguida, o Capítulo 5 apresenta o sistema implementado, suas interfaces gráficas e seus fluxos operacionais por perfil de acesso. O Capítulo 6 expõe e discute os resultados obtidos na validação empírica da plataforma junto aos usuários. Por fim, o Capítulo 7 encerra a monografia com as conclusões, a síntese do atendimento aos objetivos, as limitações identificadas e as propostas para trabalhos futuros.

\section{Justificativa}

O desenvolvimento do Test-Manager justifica-se pela necessidade prática de substituir o gerenciamento de testes descentralizado e informal no IFES Campus Alegre por uma plataforma dedicada. Atualmente, o uso de ferramentas genéricas em nuvem gera desorganização de dados e perda crônica do histórico de execuções \cite{medeiros2022, kokubum2025}. Tecnicamente, a solução proposta centraliza a gestão de \textit{Quality Assurance} (QA) e garante a rastreabilidade de defeitos por meio de um fluxo estruturado e auditável \cite{tebes2022}. No âmbito acadêmico e institucional, o projeto consolida o aprendizado prático da arquitetura MVC com o \textit{framework} Laravel \cite{subecz2021} e fornece aos coordenadores de projeto um mecanismo ágil para monitorar e avaliar a estabilidade dos sistemas desenvolvidos no \textit{campus}.

\section{Problema de Pesquisa}

O desenvolvimento de software no Instituto Federal do Espírito Santo (IFES) Campus Alegre carece de uma infraestrutura padronizada para a etapa de Garantia da Qualidade (QA), especialmente no que tange à execução de testes e ao rastreamento de defeitos. Atualmente, as equipes gerenciam seus ciclos de teste de maneira descentralizada e improvisada, recorrendo frequentemente a documentos avulsos em nuvem, como o Google Docs. 

Essa abordagem não especializada resulta em problemas crônicos no ecossistema de desenvolvimento da instituição. A literatura da área aponta que a ausência de ferramentas adequadas e a falta de alinhamento conceitual sobre os fluxos de validação geram perda constante de histórico de execuções, desorganização estrutural dos dados e extrema dificuldade na extração de métricas de qualidade confiáveis \cite{tebes2022, kokubum2025, medeiros2022}. Sem um repositório centralizado, torna-se inviável para os professores e coordenadores monitorar o progresso dos alunos e avaliar com precisão a estabilidade das aplicações desenvolvidas antes de suas respectivas entregas e implantações.

Diante da necessidade de substituir essas práticas fragmentadas por um fluxo de trabalho estruturado, seguro e auditável, a presente pesquisa é norteada pelo seguinte problema: \textbf{Como um sistema web centralizado pode melhorar o controle de testes e o registro de defeitos nos projetos de software do IFES Campus Alegre?}

\section{Objetivos}

\subsection{Objetivo Geral}
Desenvolver um sistema web centralizado para o gerenciamento de casos de teste e rastreamento de defeitos de software, denominado Test-Manager, utilizando a linguagem PHP com o \textit{framework} Laravel e o banco de dados MySQL, visando otimizar o controle de qualidade, a preservação do histórico e a governança dos projetos de software desenvolvidos no IFES Campus Alegre.

\subsection{Objetivos Específicos}
Para alcançar o objetivo geral proposto, definem-se os seguintes objetivos específicos:

\begin{itemize}
    \item Mapear e especificar detalhadamente os requisitos funcionais e não funcionais da plataforma junto às equipes de desenvolvimento e coordenação do IFES Campus Alegre;
    \item Modelar a arquitetura lógica da aplicação web e a estrutura do banco de dados relacional com base no padrão de arquitetura \textit{Model-View-Controller} (MVC);
    \item Implementar o sistema web Test-Manager contemplando os módulos de cadastro de projetos, criação e execução de casos de teste, registro analítico de defeitos e o controle hierárquico de acessos (Coordenador, Desenvolvedor e Tester);
    \item Executar testes funcionais e validações na plataforma desenvolvida para garantir a estabilidade, a segurança e a aderência do produto final às necessidades institucionais.
\end{itemize}

